\section{Load Imbalance Across Frontier Levels (Detailed)}
\subsection{Level-synchronous execution model}
Let level $L$ produce task set $\{\tau_i\}_{i=1}^{m_L}$, where each task is one chain settle for key $dk_i$.
With $P$ worker threads and a barrier at level end, level makespan is:
\[
T_L = \max_{p\in\{1,\dots,P\}} \sum_{\tau_i\in\mathcal{W}_p} C_i + B_L,
\]
where $C_i$ is chain cost and $B_L$ is barrier/synchronization overhead.

If costs are highly skewed, then $T_L$ is dominated by a small subset of chains.

\paragraph{Plain-language reading of the formula.}
At one frontier level, each thread gets a bag of chain tasks.
The level cannot finish until the \emph{slowest} thread finishes its bag.
So the term
\[
\max_p \sum_{\tau_i\in\mathcal{W}_p} C_i
\]
means ``the most overloaded thread decides the level time.''
The extra term $B_L$ is the fixed wait/synchronization cost paid by everyone at the end.

\paragraph{Why this matters.}
Even if 31 threads finish quickly, one thread with a very heavy chain still blocks progression to the next level.
This is exactly why high core count does not guarantee high speedup in level-synchronous symbolic search.

\subsection{Why successor fanout creates skew}
For a chain keyed by $dk$:
\[
C(dk) \approx c_1\,|\Pend_{dk}|\,|\Live_{dk}| + c_2\,|\Pend_{dk}|^2 + c_3\,M(dk),
\]
where:
\begin{itemize}
  \item $|\Pend_{dk}|$: new zones this level,
  \item $|\Live_{dk}|$: current antichain size,
  \item $M(dk)$: successful/attempted convex-union operations.
\end{itemize}
Uneven successor fanout makes $|\Pend_{dk}|$ heavy-tailed across keys.
If large fanout also lands on already-large antichains, cost becomes superlinear.

\paragraph{What ``uneven successor fanout'' means concretely.}
From one frontier level, each parent symbolic state generates a variable number of successors.
Those successors are then projected to keys $dk$.
If many parents happen to project to the \emph{same} key, that key's pending list explodes.
Other keys may receive only one or two pending zones.

\paragraph{Why this happens in timed symbolic models.}
This concentration is structural, not random noise:
\begin{itemize}
  \item synchronization patterns can funnel many transitions into the same discrete control state,
  \item zone abstraction merges many concrete valuations under identical discrete locations/integers,
  \item extrapolation/subsumption can keep some keys compact while others remain large and active.
\end{itemize}
So the algorithm naturally creates hot and cold keys at the same level.

\paragraph{Interpreting the proxy cost terms.}
\begin{itemize}
  \item $c_1|\Pend||\Live|$: each pending candidate is compared against many live survivors.
  \item $c_2|\Pend|^2$: pending candidates also interact indirectly through kill/merge effects.
  \item $c_3M(dk)$: successful/attempted convex unions add extra matrix work.
\end{itemize}
When both $|\Pend|$ and $|\Live|$ are large, this key becomes a straggler.

\subsection{Starvation (tail) condition}
Define total useful work in a level:
\[
W_L = \sum_{i=1}^{m_L} C_i,
\]
and ideal balanced time $W_L/P$.
A practical starvation indicator is:
\[

\paragraph{Interpretation of $\rho_L$.}
If $\rho_L=1$, the heaviest chain is about as large as ``fair share'' per worker.
If $\rho_L=4$, then one chain alone costs four workers' fair-share budget.
In that situation, no scheduling trick can fully hide the tail at that level; most workers will finish early and wait.

\paragraph{A concrete mini-scenario.}
Assume $P=8$ threads and total level work $W_L=800$ units, so ideal is $W_L/P=100$.
If one hot chain has $C_{max}=320$, then
\[
\rho_L = 320/100 = 3.2.
\]
That single chain already needs over three ideal worker budgets, so noticeable idle time is unavoidable unless the chain itself can be split.
\rho_L = \frac{\max_i C_i}{W_L/P}.
\]
When $\rho_L \gg 1$, one chain alone exceeds ideal per-thread budget and threads must idle near barrier.

\subsection{Observed effect in this checker}
The checker is level-synchronous; therefore tail chains directly delay the next level.
Even with dynamic scheduling, once only one or two heavy chains remain, available parallelism collapses.
This explains why adding threads can stop helping (or regress due to overhead).

\paragraph{Thread timeline intuition.}
At the start of a level, all workers are busy.
As light chains complete, workers become idle.
Near level end, only the hot chains remain, often on one or two workers.
The remaining workers wait at the barrier.
This ``long tail'' repeats every level where chain skew is high.

\paragraph{Key takeaway.}
The bottleneck is not ``insufficient thread count.''
The bottleneck is \emph{per-level skew in chain costs} under barrier synchronization.
So the primary optimization target is tail reduction: cost-aware ordering, heavy-chain chunking, and adaptive policy selection.

\begin{figure}[t]
\centering
\begin{tikzpicture}[x=0.19cm,y=0.62cm,>=Latex, font=\small]
  \foreach \i/\lbl in {0/T0,1/T1,2/T2,3/T3,4/T4,5/T5,6/T6,7/T7} {
    \draw[gray!55] (0,-\i) -- (42,-\i);
    \node[left] at (0,-\i) {\lbl};
  }

  % light tasks
  \filldraw[draw=black,fill=blue!20] (0,-0.3) rectangle (8,0.3);
  \filldraw[draw=black,fill=blue!20] (0,-1.3) rectangle (7,-0.7);
  \filldraw[draw=black,fill=blue!20] (0,-2.3) rectangle (9,-1.7);
  \filldraw[draw=black,fill=blue!20] (0,-3.3) rectangle (6,-2.7);
  \filldraw[draw=black,fill=blue!20] (0,-4.3) rectangle (8,-3.7);
  \filldraw[draw=black,fill=blue!20] (0,-5.3) rectangle (7,-4.7);
  \filldraw[draw=black,fill=blue!20] (0,-6.3) rectangle (9,-5.7);

  % one heavy chain
  \filldraw[draw=black,fill=red!25] (0,-7.3) rectangle (34,-6.7);

  % idle tails
  \foreach \i in {0,1,2,3,4,5,6} {
    \filldraw[draw=black,fill=gray!20] (9,-\i-0.3) rectangle (34,-\i+0.3);
  }

  \draw[dashed] (34,0.8) -- (34,-7.8);
  \node[above] at (34,0.8) {level barrier};

  \node at (17,-7.0) {hot chain};
  \node at (21,-1.0) {idle waiting};
\end{tikzpicture}
\caption{Starvation/tail effect: one hot chain dominates level completion while other threads idle.}
\label{fig:parallel-starvation}
\end{figure}

